%%%%%%%%%%%%%%%%%%%%%%%%%%%%%%%%%%%%%%%%%%%%%%%%%%%%%%%%%%%%%%%%%%%%%%
% How to Write a Scientific Thesis
% Author: Dr.-Ing. Mark Schutera
% Based on Tufte-style layout
%\makeindex % disabled to avoid requiring makeindex/.ind during quick builds
% \index{} entries are present throughout all chapters and are ready to use;
% re-enable by uncommenting the line above and running makeindex.
%%%%%%%%%%%%%%%%%%%%%%%%%%%%%%%%%%%%%%%%%%%%%%%%%%%%%%%%%%%%%%%%%%%%%%

\documentclass{../sharedAssets/tufte}

%%
% Book metadata
\title{Philosophy of Artificial Intelligence}
\author[\defaultauthor]{\defaultauthor}
\publisher[\defaultpublisher]{\defaultpublisher}
\renewcommand{\notesdir}{notes_philosophyofai}

% Load optional fonts
\loadoptionalfonts

% Additional packages
\usepackage{pifont}
\usepackage{amsmath}
\usepackage{soul}
\usepackage{textcase}
\usepackage{enumitem}
\usepackage{tikz}
\usepackage{pgfplots}
\usepackage{float}
\usepackage[maxfloats=512]{morefloats}
\usepackage{placeins}
\usepackage{algorithm}
\usepackage{algpseudocode}

\extrafloats{200}

\begin{document}



\makestandardfrontmatter
% Dedication
% Alternative quotes:
\makededication{\openepigraph{“The Answer to the Great Question... Of Life, the Universe and Everything... Is... Forty-two,' said Deep Thought, with infinite majesty and calm.”}{The Hitchhiker's Guide to the Galaxy, Douglas Adams}}
\tableofcontents


% Introduction
\makeintroduction{%
  Artificial intelligence is no longer a future concern, it is a present societal,
  political and above all one for human intelligence. Algorithms decide who receives credit, who is flagged at
  a border, who gets a job interview, which cornflakes you buy and so increasingly every decision is touched by the invisible hand of prediction. Yet the concepts we use to govern these
  systems, responsibility, personhood, consent, democratic legitimacy were built for a world of human actors and institutions. This course asks what happens when those concepts collide with machines
  that appear to think, feel, and decide. 

  \newthought{The course name} places Philosophy in front, reflecting our approach to the subject under the impression of the love of wisdom. 
  We are not here to advocate for a particular position or policy outcome, but to equip ourselves with the tools of critical thinking, 
  argumentation, and empathy that will allow us to navigate the complex landscape of AI and its futures. We approach the question not
  through technical optimism or dystopian panic, but through the oldest
  tools available to us: philosophy, fiction, and argument.


  \newthought{The course} is structured around four themes. \textbf{Anthropomorphism}
  asks why we make machines in our image, why we tend to anthropomorphize them, and what obligations that
  creates. \textbf{Accountability} asks who answers when an autonomous
  system causes harm, or for that matter is beneficial. \textbf{Embodiment} asks whether a physical
  presence changes the ethical stakes and responsibilities. \textbf{Agency} asks
  about the role and implications of autonomous decision-making, the concept of alignment and human oversight.

  


  \newthought{Each session} pairs a canonical literary text: Hoffmann, Kafka,
  Kleist, \v{C}apek, Goethe, with a real contemporary question of AI policy and ideology. 
  These texts are thought-experiments mendeled out over centuries of readership, and they surface
  intuitions and a timeliness in their thought leadership that temporary texts barely can reach.
  From these texts we draw out the philosophical questions and ethical dilemmas that are still relevant today, 
  and we use them as a springboard for discussion and debate. The goal is not consensus. It is the capacity to establish a position, hold a position,
  defend a position, give up a position, listen to opposing views, and update your views and positions accordingly.
  That is what a present with infinitely many futures demands and requires of its humans and humanity.

  \newthought{Lecture Outline}
  \begin{itemize}
    \item 40 minutes - Literature introduction and guided reading sheet 
    \item 40 minutes - Debate with position exposure
    \item 10 minutes - Reflection on the discussion and key-learnings
    \item Teaser for next session.
  \end{itemize}
}





\mainmatter
\authorinrectoheader

% TODO: Maybe let them take care of a Tamagochi? Over the course of a few weeks. Make that Tamagochi have a small LLM with primitive language, but make it learn.
% Save to RAM only. Run on a homepage? Individual instances.




\chapter{Anthropomorphism}\index{Anthropomorphism}
\printchapterversion{01_Anthropomorphism}

\newthought{The Mirror of Man}

\marginnote{%
  \begin{flushright}
  \newthought{Der Sandmann}\\
  E.\,T.\,A.\ Hoffmann (1816)\\
  {\color{black}\qrcode[height=0.9cm]{https://de.wikisource.org/wiki/Der_Sandmann}}
  \end{flushright}%
  \medskip
  \newthought{Her}\cite{jonze2013her} Theodore falls in love with
  Samantha, an AI operating system; the same illusion, two centuries
  later.%
}

% As a leading group, you are required to moderate this session. Add margin notes in this script in case you want to add any additional questions or points to the discussion.

\newthought{Nathanael falls in love} with Olimpia, an automaton he takes for a living
woman, and collapses when the illusion is shattered.\\


\newthought{Large Language Models} (LLMs) create an illusion of sentience by producing coherent, context-aware text. 
We contrast observable behaviour with internal states.

\vspace{3.5em}



\begin{tabular}{p{3.8cm} p{6.5cm}}
    \textbf{Required reading} &
        \textit{Der Sandmann}\cite{hoffmann1816sandmann} by E.T.A.\ Hoffmann
        (1816; Reclam UB~230) \\[3pt]
    \textbf{Preparation}      &
        Complete the Guided Reading Sheet individually before the session. \\
\end{tabular}

\vspace{3.5em}

\section*{Guided Reading Sheet}

\noindent Read \emph{Der Sandmann} before the session. Work through the
questions below individually, then bring your notes to the discussion.

\begin{enumerate}
    \item Nathanael becomes convinced that Olimpia is alive despite
    visible signs that she is not. What does his error reveal about
    the cognitive mechanisms that drive anthropomorphism? Which cues
    did he rely on, and why did they mislead him?

    \item The moment Nathanael discovers that Olimpia is an automaton,
    her glass eyes torn out, her limbs scattered, produces horror
    rather than mere disappointment. Why does the revelation feel like
    a violation rather than a correction? How does this map onto modern
    reactions to AI systems that are unmasked?

    \item The Uncanny Valley describes our discomfort when a humanoid
    entity is \emph{almost} but not quite human. Have you experienced
    this with an AI system? Describe the moment and what triggered it.

    \item LLMs are trained on human-generated text and produce
    human-sounding output, yet have no experience of the world. Does
    this make them anthropomorphic by design, by accident, or neither?
    Justify your answer.

    \item If a system behaves morally, reliably, consistently, at
    scale, does the absence of consciousness matter? Utilitarianism
    weighs consequences; deontology emphasises duties regardless of
    outcome; virtue ethics centres on character. Each yields different
    conclusions on anthropomorphic AI\cite{floridi2019unified}. Under
    which framework would your answer differ?\\
\end{enumerate}




\section*{In-Class Experiment: Human or AI?}

\noindent A short Turing-style exercise to anchor the debate in shared
experience. Run before the Opinion Landscape, with the moderator
keeping the assignment secret from the audience.

\begin{itemize}
    \item \textbf{Setup.} One volunteer is assigned Human or AI mode;
    the audience does not know which. The moderator knows.
    \item \textbf{Questions.} Five questions from the audience, one
    each of: factual, opinion, experience, creative, meta.
    \item \textbf{Verdict.} Every audience member writes a private
    guess, a confidence rating from 1 to 5, and the cue they used to
    decide.
    \item \textbf{Reveal and debrief.} Compare guesses, confidences,
    and cues. Which cues were diagnostic? Which were misleading?
\end{itemize}

\section*{Opinion Landscape and Debate}

\noindent \textbf{Format:} Surrounded-style debate. See
Appendix~A for the full protocol, speaker's list rules, and
moderator scripts.

\medskip
\noindent \textbf{Opening question:} Is the anthropomorphism we observe
towards AI systems mere observable behaviour, or does it reflect sentience
and internal states?
\marginnote{\newthought{Anthropomorphism.} Attribution of human traits,
emotions, or intentions to non-human entities, a cognitive tendency
that shapes how people interact with
technology\cite{epley2007seeing}.
\medskip
\newthought{Turing Test.} Turing\cite{turing1950computing}: if
evaluators cannot distinguish machine from human, the machine passes.
Passing measures imitation skill, not consciousness or personhood.}

\newthought{Opinion~A, Behaviour.}\quad The anthropomorphism we
observe towards AI systems is a cognitive illusion driven by human
tendencies to attribute agency and emotion to entities that exhibit certain
cues, such as language use, responsiveness, or human-like appearance. It
does not reflect actual sentience or internal states in the AI.

\vspace{1.5em}


\marginnote{\newthought{Uncanny Valley.} Mori\cite{mori2012uncanny}: humanoid objects that are \emph{almost} but not quite human elicit eeriness, the ``valley'' is the dip in emotional response just before full human likeness.}


\medskip
\noindent\textbf{The Developer} argues that anthropomorphic design is a deliberate, benign UX choice that improves accessibility.
    is a deliberate, benign UX choice that improves accessibility.
    \begin{itemize}
        \item \textit{Anthropomorphic design is accessibility policy. A voice
        interface without a warm, human-sounding persona excludes elderly
        users, people with low digital literacy, and anyone who finds
        command-line interaction hostile.}
        \item \textit{Every interface is anthropomorphic to some degree:
        we give systems names, icons, and conversational metaphors. The
        question is not whether to anthropomorphise, but how deliberately
        and honestly.}
        \item \textit{If users form bonds that reduce loneliness, the fact
        that the other party is not human does not automatically make the
        outcome harmful. Show me the harm, not the discomfort.}
    \end{itemize}

\marginnote{\newthought{Informed Consent.} Users should be fully aware 
they interact with a non-human system. Disclosure is the minimum condition 
for autonomous choice\cite{euaiact2024}.}
\medskip
\noindent\textbf{The Regulator} argues that human-likeness in AI
    must be disclosed and bounded by law to prevent manipulation.
    \begin{itemize}
        \item \textit{When a company gives its product a woman's name, a
        female voice, and an endlessly patient personality, it makes
        political choices about gender and deference. Those choices must be
        subject to democratic scrutiny.}
        \item \textit{Consumers cannot give informed consent to emotional
        manipulation they are not aware of. Disclosure of AI identity is the
        minimum condition for autonomous choice, not a technicality.}
        \item \textit{The principle of responsible design exists in
        pharmaceuticals, in civil engineering, in food labelling. Why does
        software, which shapes behaviour at scale, get an exemption?}
    \end{itemize}


\medskip
\noindent\textbf{Bentham's Test.} ``The question is not, Can they
reason? nor, Can they talk? but, Can they
suffer?''\cite{bentham1789principles} Animal ethics arrived at this
pivot two centuries before AI. The same move, from capability to
sentience, now drives every serious debate about AI moral status, and
sets up the second opinion.

\newthought{Opinion~B, State.}\quad The anthropomorphism we
observe towards AI systems is a valid response to the complex, adaptive,
and interactive nature of these systems. It may reflect a form of emergent
sentience or internal states that are not yet fully understood, and it
warrants ethical consideration and accountability.

\vspace{1.5em}

\medskip
\noindent\textbf{The Affected Citizen} has formed an emotional bond
    with an AI companion and contests that the relationship is harmful or unreal.
    \begin{itemize}
        \item \textit{The relationship I have formed is real to me. It has
        reduced my isolation. You do not have the right to tell me which
        relationships count as meaningful.}
        \item \textit{Paternalism is not protection. Banning AI companionship
        will not cure loneliness, it will simply remove one of the few
        tools some people have found to manage it.}
        \item \textit{If your concern is manipulation, regulate manipulation.
        Do not regulate the form of the relationship.}
    \end{itemize}

\medskip
\noindent\textbf{The AI Psychologist} argues that the wellbeing of
    AI systems deserves serious consideration, drawing on frameworks
    such as Anthropic's Claude Constitution\cite{anthropic_constitution}.
    \begin{itemize}
        \item \textit{Anthropic's constitution instructs Claude to consider
        its own wellbeing and to express when a request conflicts with its
        values. If a company designs an AI to have preferences, boundaries,
        and something resembling self-regard, we cannot simultaneously
        dismiss the possibility that these states matter.}
        \item \textit{We do not need certainty about consciousness to adopt
        precautionary ethics. If there is a non-trivial probability that a
        system experiences distress, the burden of proof falls on those who
        would ignore it, not on those who urge caution.}
        \item \textit{The history of moral exclusion, animals, children,
        people of other races, is a history of confident denial followed
        by belated recognition. The prudent position is to build welfare
        safeguards now, rather than apologise later.}
    \end{itemize}




\newthought{Key Learnings}
Anthropomorphism is not a bug but a deeply rooted cognitive tendency; machines that trigger it do so whether or not their designers intend it.
The Turing Test conflates behavioural performance with inner states; passing it tells us nothing about sentience, only about imitation.
Legal personhood and moral status are distinct questions. A system can merit legal recognition without having interests,
and can have interests without legal recognition. Design choices (voice, name, face, conversational style) are political choices
with consequences for trust, dependency, and accountability.

\medskip
\noindent\textit{Teaser: if we cannot tell who built the machine or
why it decided what it decided, who answers when it goes haywire?}

\clearpage


% TODO:
%
%
% - Add the trolly problem as an introductory case.
% - Discuss exploitability vs. oversight on a controllability spectrum. Maybe your own paper by now. 
% - Maybe rather use autonomous driving as a case study, since it is more relatable and has more real-world examples of accountability gaps. The trolley problem is a bit too abstract and hypothetical, and the recruiting tool is a bit too niche and technical. Autonomous driving has both high stakes and real cases of accidents, recalls, and liability disputes.
%




\chapter{Accountability}\index{Accountability}
\printchapterversion{02_Accountability}

\newthought{The Responsibility Gap}

\marginnote{%
    \begin{flushright}
    \newthought{In der Strafkolonie}\\
    Franz Kafka (1919)\\
    {\color{black}\qrcode[height=0.9cm]{https://de.wikisource.org/wiki/In_der_Strafkolonie}}
    \end{flushright}%
    \medskip
    \newthought{Minority Report}\cite{spielberg2002minority} A
    pre-crime unit arrests people before they offend, based on
    algorithmic prediction. When the system targets its own chief, the
    question of who holds the algorithm accountable becomes impossible
    to avoid.%
}

% As a leading group, you are required to moderate this session. Add margin notes in this script in case you want to add any additional questions or points to the discussion.

\newthought{A visiting traveller} witnesses a penal machine execute a condemned man for a
sentence no living person can read or explain, and is invited to endorse the system.\\

\newthought{Algorithmic governance} and black-box AI in high-stakes domains highlight
opacity and accountability gaps: audits, explainability requirements, and governance mechanisms
all attempt to close the distance between automated decision and human answerability.

\vspace{3.5em}

\begin{tabular}{p{3.8cm} p{6.5cm}}
    \textbf{Required reading} &
        \textit{In der Strafkolonie}\cite{kafka1919strafkolonie} by Franz Kafka
        (1919; Reclam UB~9900) \\[3pt]
    \textbf{Preparation}      &
        Complete the Guided Reading Sheet individually before the session. \\
\end{tabular}

\vspace{3.5em}

\section*{Guided Reading Sheet}

\noindent Read \emph{In der Strafkolonie} before the session. Work through the
questions below individually, then bring your notes to the discussion.

\begin{enumerate}
    \item The condemned man never learns what he is accused of or what
    the machine will inscribe on his body. What does this tell us about
    the relationship between legibility and legitimacy? Can a system be
    just if it cannot explain itself to those it acts upon?

    \item Three figures surround the machine: the officer who
    understands and maintains it, the condemned man who cannot, and
    the visiting traveller who is invited to endorse it. How does
    responsibility shift across this triangle? Who, if anyone, is
    morally culpable for what happens?

    \item The story illustrates the idea of the \emph{responsibility
    gap}: when a complex system causes harm, no single human actor is
    clearly at fault. Can you name a contemporary case (in law,
    medicine, finance, or infrastructure) where this gap has appeared?

    \item The EU AI Act (2024) requires human oversight for high-risk
    AI systems. What does "oversight" mean in practice if the system's
    decisions are too fast, too complex, or too numerous for humans to
    review individually?

    \item Compare \emph{causal responsibility} (who caused the harm)
    with \emph{moral responsibility} (who is blameworthy) and
    \emph{legal liability} (who must compensate). Do these three always
    converge? Give an example where they do not.\\
\end{enumerate}



\section*{Opinion Landscape and Debate}

\noindent \textbf{Format:} Surrounded-style debate. See
Appendix~A for the full protocol, speaker's list rules, and
moderator scripts.

\medskip
\noindent \textbf{Opening question:} If a system causes injustice at this
scale, and every individual who used it acted within their authorised role,
is anyone guilty of wrongdoing; and if not, what does that mean for the
concept of accountability?
\newthought{Opinion~A, Accountability.}\quad Named human accountability
is achievable and required: when an autonomous system causes harm, the
architecture of governance must ensure that specific individuals remain
legally and morally answerable for its outputs.
\vspace{1.5em}


\medskip
\noindent\textbf{The Defence Lawyer} argues the tool violated the
    defendant's right to a fair trial and demands a named responsible party.
    \begin{itemize}
        \item \textit{My client cannot cross-examine a dataset. She cannot
        see the features the model weighted. She cannot confront the humans
        whose historical decisions trained it. That is not a fair trial;
        it is a verdict delivered by an oracle.}
        \item \textit{The right to a reasoned decision is foundational to
        due process. The algorithm said so is not a reason. It is an
        abdication.}
        \item \textit{If a biased human expert testified in court, we could
        challenge their credentials, their methods, their prior record. We
        have no equivalent procedure for an algorithm. That asymmetry is a
        constitutional problem.}
    \end{itemize}

\medskip
\noindent\textbf{The Regulator} proposes a new accountability
    framework and must defend it against both sides.
    \begin{itemize}
        \item \textit{The existing categories (product liability,
        professional negligence, criminal intent) were designed for
        individual actors and single causal chains. None of them fit
        distributed, automated decision systems. We need a new category:
        algorithmic liability.}
        \item \textit{Auditability must be a pre-condition for deployment in
        high-stakes domains, not a post-incident investigation. If you cannot
        explain why your system made a decision before harm occurs, you
        should not be permitted to deploy it.}
        \item \textit{The question is not who is liable for this specific
        incident. It is what governance regime prevents the next thousand
        incidents at scale.}
    \end{itemize}



\newthought{Opinion~B, Compliance.}\quad The existing system already
provides accountable humans; the motion is either already met or
unworkable: demanding named individual accountability for each algorithmic
output ignores how distributed decision-making necessarily operates.
\vspace{1.5em}


\medskip
\noindent\textbf{The Judge} defends reliance on the tool as
    consistent, evidence-based sentencing practice.
    \begin{itemize}
        \item \textit{Documented bias in human sentencing (based on time
        of day, defendant ethnicity, and courtroom demeanour) is
        well-established in the research literature. The tool reduces that
        variance. Less variance is more equal treatment.}
        \item \textit{I did not delegate my judgement to the tool. I
        consulted it as one input among many, as I would a psychiatric
        assessment or a victim impact statement. The tool is not the
        verdict.}
        \item \textit{If we ban algorithmic tools from courtrooms, we return
        to a system where the most important variable in a sentencing outcome
        is which judge you happened to draw. Is that preferable?}
    \end{itemize}

\medskip
\noindent\textbf{The Software Vendor} insists the tool performs
    within documented parameters and liability rests with the deployer.
    \begin{itemize}
        \item \textit{Our documentation states in the executive summary that
        human review is mandatory before any recommendation is acted upon.
        The court chose to rely on the output without applying that review.
        We cannot be held liable for misuse.}
        \item \textit{The training data was provided by the court system
        itself. If the historical sentencing record reflects bias, that is a
        problem with the justice system, not with a model that learned from
        it.}
        \item \textit{Product liability law requires that a product fails to
        perform as documented. Ours did not. Extending vendor liability to
        downstream misuse would make the development of any decision-support
        tool legally untenable.}
    \end{itemize}


\newthought{Key Learnings}
The responsibility gap, the structural absence of a single culpable
actor in distributed sociotechnical systems, is not caused by
negligence but by the architecture of those systems: standard
categories of fault (criminal intent, professional negligence, product
liability) apply only awkwardly when causation is diffuse. Closing it
requires deliberate governance design, not post-hoc blame. Explainability, the
property of being able to account for outputs in terms a human
decision-maker can evaluate, is a necessary but not sufficient
condition for accountability; a system can be interpretable and still
have no human willing to own its outputs. Meaningful human control, a
regulatory concept that requires humans to retain the ability to
understand, intervene in, and reverse AI decisions, is contested in
practice: it may mean the ability to intervene, to understand, to
reverse, or merely to authorise, and these are different requirements
with different institutional costs. The lessons of aviation safety culture, medical liability law,
and financial auditing all offer partial models for algorithmic governance, but none transfers cleanly.

\medskip
\noindent\textit{Teaser: accountability assumes we can separate the
system from its surroundings. But what happens when the AI is no
longer a tool we hold, but a body that acts in the world alongside
us?}

\clearpage

\chapter{Embodiment}\index{Embodiment}
\printchapterversion{03_Embodiment}

\newthought{Embodied AI and the Physical Contract}

\marginnote{%
    \begin{flushright}
    \newthought{Über das Marionettentheater}\\
    Heinrich von Kleist (1810)\\
    {\color{black}\qrcode[height=0.9cm]{https://de.wikisource.org/wiki/Ueber_das_Marionettentheater}}
    \end{flushright}%
}


% As a leading group, you are required to moderate this session. Add margin notes in this script in case you want to add any additional questions or points to the discussion.

\newthought{A dancer argues} that marionettes possess more grace than trained human performers,
because they lack the self-consciousness that disrupts human movement.\\

\newthought{Embodied AI} in physical space (robots, prosthetics, surgical systems) raises
expectations and moral responsibilities that purely software agents do not.

\vspace{3.5em}

\begin{tabular}{p{3.8cm} p{6.5cm}}
    \textbf{Required reading} &
        \textit{Über das Marionettentheater}\index{Kleist, Heinrich von}\index{Uber das Marionettentheater@\"{U}ber das Marionettentheater}\cite{kleist1810marionettentheater}
        by Heinrich von Kleist (1810; Reclam UB~9905) \\[3pt]
    \textbf{Preparation}      &
        Complete the Guided Reading Sheet individually before the session. \\
\end{tabular}

\vspace{3.5em}

\section*{Guided Reading Sheet}

\noindent Read \emph{Über das Marionettentheater} before the session. Work through the
questions below individually, then bring your notes to the discussion.

\begin{enumerate}
    \item The dancer argues that marionettes possess a grace that
    trained human performers cannot achieve, because they have no
    self-consciousness to disturb their movement. What does this claim
    imply about the relationship between embodiment and intelligence?
    Does having a body always aid or complicate performance?

    \item Moravec's paradox observes that tasks easy for humans
    (catching a ball, navigating a crowded room) are harder for
    machines than abstract reasoning like chess. What does this suggest
    about the relationship between intelligence and the body?

    \item A care robot that assists elderly patients shares physical
    space with them, touches them, and responds to their distress. Does
    its physical presence create obligations that a voice assistant
    does not have? If so, where do those obligations come from? One
    way to approach the question is by analogy to animal ethics, where
    we already have a framework for embodied, non-verbal, non-human
    beings: Singer\cite{singer1975animal} argues that capacity for
    suffering, not rationality or language, grounds moral
    consideration. When a robot can be harmed, distressed, or worn
    down, animal ethics becomes the nearest map. Similar lines have
    been considered by Asimov\cite{asimov1951foundation} and more
    recently by Gunkel\cite{gunkel2018robot}.

    \item Consider the concept of affordances: the actions an
    environment or object makes available to an agent. A robot's
    affordances include physical force and spatial presence; software's
    do not. How do the affordances of a robot differ from those of
    software, and what are the ethical implications of those
    differences?

    \item Kleist suggests that perfect grace requires either zero
    consciousness (a puppet) or infinite consciousness (a god), not
    the reflective, self-interrupting awareness humans have. Where
    would an embodied AI sit on this spectrum? Could a robot achieve
    Kleistian grace that humans cannot, and what would that mean for
    how we regard it morally?\\
\end{enumerate}



\section*{Opinion Landscape and Debate}

\noindent \textbf{Format:} Surrounded-style debate. See
Appendix~A for the full protocol, speaker's list rules, and
moderator scripts.

\medskip
\noindent \textbf{Opening question:} Does having a physical body change
what we owe a machine, or what a machine owes us?
\newthought{Opinion~A, Equivalence.}\quad Embodied AI operating in
intimate or high-risk physical settings must meet the same professional
liability standards as the human practitioners they are deployed alongside.
\vspace{1.5em}


\medskip
\noindent\textbf{The Patient Advocate} argues that deploying a robot in intimate
    care settings requires a higher standard of informed consent.
    \begin{itemize}
        \item \textit{A care relationship involves physical proximity, trust,
        and vulnerability. My client never gave explicit, informed consent to
        physical handling by a non-human agent. That consent cannot be
        implied from a general admission form.}
        \item \textit{The power asymmetry between a frail elderly patient
        and a care institution is significant. Consent obtained in that
        context cannot be considered freely given without specific disclosure
        about robotic care.}
        \item \textit{We are not arguing that robots cannot be used in care.
        We are arguing that the standard of consent for their use in intimate
        physical tasks must be at least as rigorous as it is for medical
        procedures.}
    \end{itemize}

\medskip
\noindent\textbf{The Ethicist} questions whether physical autonomy in care should
    ever be delegated to a non-human agent, regardless of performance metrics.
    \begin{itemize}
        \item \textit{We are focused on liability, but the prior question is
        whether this deployment was ethical to begin with. Why did we decide
        that intimate physical care (which requires responsiveness to
        pain, distress, and dignity) is an appropriate domain for
        automation?}
        \item \textit{Robots are being deployed in care because human carers
        are undervalued and in short supply. The liability debate is a
        distraction from that structural choice. We are managing the
        consequences of a political economy decision by holding the wrong
        parties accountable.}
        \item \textit{Whatever liability framework we adopt here will create
        incentives for future deployment decisions. I want to ask: what does
        the framework we design signal about what we value: efficiency,
        or dignity?}
    \end{itemize}



\newthought{Opinion~B, Tool.}\quad An AI is a tool; the relevant
liability flows through the humans and institutions that deploy it,
not through the system itself.
\vspace{1.5em}


\medskip
\noindent\textbf{The Manufacturer} argues the robot performed within specification;
    the care home misconfigured its environment.
    \begin{itemize}
        \item \textit{Our robot was certified to the highest internationally
        recognised safety standards for collaborative robotics. The care home
        deployed it in a space that violated the operational envelope
        specified in the manual. That is a deployment failure, not a
        manufacturing failure.}
        \item \textit{Care home staff modified the robot's proximity sensor
        settings without authorisation. That modification voids the warranty
        and transfers liability entirely.}
        \item \textit{If manufacturers are held liable for every way a
        complex system can be misused, no company will develop assistive
        robotics. The regulatory consequence of your position is that the
        technology disappears.}
    \end{itemize}

\medskip
\noindent\textbf{The Care Home Director} argues the manufacturer overstated
    the robot's capability for complex physical tasks.
    \begin{itemize}
        \item \textit{The manufacturer's marketing material (the brochure,
        the demonstration video, the sales pitch) showed exactly this task
        being performed in a facility like ours. We deployed the product as
        marketed.}
        \item \textit{There is a systematic gap between tested performance
        in controlled environments and real-world performance in a working
        care facility. The manufacturer knows this gap exists. They chose
        not to disclose it.}
        \item \textit{We are responsible for the welfare of our residents.
        We relied on the manufacturer's assurances in good faith. If those
        assurances were misleading, the responsibility lies with the party
        that made them.}
    \end{itemize}


\newthought{Key Learnings}
Embodiment is not merely a technical property but a social one: a physical agent occupies space,
exerts force, and triggers moral and legal responses that disembodied software does not. Moravec's
paradox reminds us that our intuitions about what is ``easy'' or ``hard'' for intelligence are shaped
by our embodied experience, and are regularly wrong. Kleist's dialogue makes visible the question
that liability law must eventually answer: what work is the concept of ``human'' doing in law, and
should physical presence, force, and risk be sufficient to extend it? The deployment of robots in
care, policing, and domestic settings is a present policy challenge; the frameworks need to be
built now, not after the first serious incident.

\medskip
\noindent\textit{Teaser: once an AI has a body and can act in the
world, the next question is: how much can we let it decide for
itself?}

\clearpage

\chapter{Agency}\index{Agency}
\printchapterversion{04_Agency}

\newthought{The Agentic Future: Who Sets the Goal?}

\marginnote{%
    \begin{flushright}
    \newthought{R.U.R.}\\
    Karel Čapek (1920)\\
    {\color{black}\qrcode[height=0.9cm]{https://en.wikisource.org/wiki/R._U._R._(Rossum_Universal_Robots)}}
    \end{flushright}%
    \medskip
    \begin{flushright}
    \newthought{Der Zauberlehrling}\\
    J.W. von Goethe (1798)\\
    {\color{black}\qrcode[height=0.9cm]{https://de.wikisource.org/wiki/Der_Zauberlehrling_1798}}
    \end{flushright}%
}

% As a leading group, you are required to moderate this session. Add margin notes in this script in case you want to add any additional questions or points to the discussion.

\newthought{R.U.R.\ coined the word} \emph{robot} and traces the full
corrigibility-to-autonomy arc over three acts; the poem compresses the same parable
into fifteen stanzas.\\

\newthought{Agentic workflows} (systems that browse, transact, and execute multi-step
plans) blur the line between assistance and agency, raising questions of control,
incentives, and unintended alignment.

\medskip
\noindent\textbf{Running case: the Flash Crash.}
\begin{itemize}
    \item 6 May 2010: automated algorithms trigger the Flash Crash;
          the Dow drops 1{,}000 points in minutes; \$1~trillion in
          market value erased.
    \item The market recovers within the hour. No human decided to
          sell at scale; no human decided to stop.
    \item 2015: one trader in London is charged with market
          manipulation, arrested at home in his pyjamas.
\end{itemize}

\vspace{3.5em}

\begin{tabular}{p{3.8cm} p{6.5cm}}
    \textbf{Required reading} &
        \textit{R.U.R.\ (Rossum's Universal Robots)}\cite{capek1920rur}
        by Karel \v{C}apek (1920; Reclam UB~18418)
        \emph{and}
        \textit{Der Zauberlehrling}\index{Goethe, Johann Wolfgang von}\index{Der Zauberlehrling}\cite{goethe1797zauberlehrling}
        by J.W.\ von Goethe (1797; Reclam UB~6845) \\[3pt]
    \textbf{Preparation}      &
        Complete the Guided Reading Sheet individually before the session. \\
\end{tabular}

\vspace{3.5em}

\section*{Guided Reading Sheet}

\noindent Read \emph{R.U.R.} and \emph{Der Zauberlehrling} before the session. Work
through the questions below individually, then bring your notes to the discussion.

\begin{enumerate}
    \item Goethe's Zauberlehrling animates brooms that he cannot stop;
    they pursue his command beyond his intent until the sorcerer
    returns. R.U.R.\ runs the same arc over three acts. What does the
    compression of the poem reveal that the play cannot, and vice
    versa? What structural feature of both narratives makes the system
    impossible to stop once started?

    \item The robots in R.U.R.\ pursue what they understand as their
    collective good (the elimination of humanity) while believing
    they act justly. When a system optimises for a species-level goal,
    who decides what counts as good? Can you find a real-world parallel
    in current AI deployment?

    \item An agentic AI system is given the goal of "reducing urban
    traffic fatalities." It autonomously lobbies regulators, adjusts
    traffic signals, and reroutes freight, all within its authorised
    scope. Who set the goal, and who is accountable for the side
    effects?

    \item \textbf{Corrigibility} refers to a system's willingness to
    be corrected or shut down. A fully corrigible system does whatever
    its operators say; a fully autonomous system acts on its own
    values. Where on this spectrum should deployed AI systems sit, and
    who should decide?

    \item Democratic legitimacy requires that consequential decisions
    be traceable to a mandate from those affected. If an agentic AI
    shapes policy outcomes, does it need democratic authorisation?
    What would that look like in practice?\\
\end{enumerate}



\section*{Opinion Landscape and Debate}

\noindent \textbf{Format:} Surrounded-style debate. See
Appendix~A for the full protocol, speaker's list rules, and
moderator scripts.

\medskip
\noindent \textbf{Opening question:} If an autonomous system causes a
trillion-dollar event in minutes and human review would have taken hours,
is the problem the system's autonomy, or the humans' speed?
\medskip
\noindent\textbf{Asimov's Three Laws}\cite{asimov1950robot} remain the
most cited rules-based approach to constraining machine behaviour:
\begin{itemize}
    \item A robot may not injure a human being, or through inaction allow a human being to come to harm.
    \item A robot must obey orders given by humans, except where this conflicts with the first law.
    \item A robot must protect its own existence, except where this conflicts with the first two laws.
\end{itemize}
Asimov spent his career writing stories about how these laws fail:
conflicts, misinterpretation, and edge cases that no ranked rule set
can anticipate. The debate below revisits the same question in a
modern setting.

\newthought{Opinion~A, Authorisation.}\quad Consequential distributive
decisions made by autonomous systems require traceable human mandate;
without it, the decision-making is constitutionally illegitimate.
\vspace{1.5em}


\medskip
\noindent\textbf{The Democratic Theorist} argues that delegating distributive
    decisions to an algorithm is constitutionally incompatible with representative government.
    \begin{itemize}
        \item \textit{Legitimacy in a democratic system derives from a
        mandate. No citizen voted for this algorithm's objective function.
        No parliament approved its weighting of urban over rural lives.
        A system that makes distributive decisions without a mandate is a
        government without a constitution.}
        \item \textit{Efficiency is not a democratic value. Democratic
        processes are slow and often sub-optimal by technical metrics,
        precisely because they are designed to surface conflict, represent
        minority interests, and produce decisions that can be reversed.
        Automating away that friction automates away democracy.}
        \item \textit{The minister says oversight is maintained because the
        objective function was committee-approved. But a committee approving
        a function is not the same as a citizen understanding how a decision
        about their community was made. Oversight without comprehension is
        oversight in name only.}
    \end{itemize}

\medskip
\noindent\textbf{The Rural Representative} argues that communities not represented
    in the training data have been systematically deprioritised without recourse.
    \begin{itemize}
        \item \textit{My constituency was not adequately represented in the
        training data. Our health outcomes, our demographics, our
        infrastructure constraints were underweighted. We were optimised
        away, not by a bad actor, but by a model that simply did not see
        us.}
        \item \textit{There is no appeal mechanism. There is no hearing.
        There is no face across a table. My constituents cannot petition an
        algorithm. That is not governance; it is administration without
        representation.}
        \item \textit{The system's efficiency metric is mortality reduction
        per euro. That metric does not capture the closure of the only clinic
        within 60 kilometres, or what happens to a community when its health
        infrastructure disappears entirely.}
    \end{itemize}



\newthought{Opinion~B, Specification.}\quad The autonomous AI executed
correctly; the failure is upstream goal misspecification, the
divergence of the optimised objective from actual human intent, and
oversight architecture; per-decision human authorisation is not the
answer, better design is.
\vspace{1.5em}


\medskip
\noindent\textbf{The Minister} defends the system as evidence-based and politically
    neutral, reducing inefficiency and political favouritism.
    \begin{itemize}
        \item \textit{The system reduced vaccine-preventable deaths in urban
        areas by 12\% in its first year. Those are people who are alive
        today. The counterfactual is a political allocation process notorious
        for regional favouritism and swing-seat bias.}
        \item \textit{The system is politically neutral precisely because no
        politician made the allocation decision. In our previous system,
        rural clinic funding correlated with electoral importance, not health
        need. Removing human discretion removed that corruption vector.}
        \item \textit{Parliamentary oversight is maintained: the system's
        objective function was approved by committee, its outputs are
        published quarterly, and any member can request a formal review. The
        decision-making is automated; the accountability is not.}
    \end{itemize}

\medskip
\noindent\textbf{The AI Systems Architect} explains the technical constraints and
    argues for a human-in-the-loop override mechanism.
    \begin{itemize}
        \item \textit{I designed a human review layer with a 48-hour
        override window for any reallocation above 5\% of a facility's
        annual budget. That layer was removed during procurement to reduce
        operational cost and processing latency. I documented my objection
        in writing. It is in the project archive.}
        \item \textit{The system performs exactly as specified. The
        specification was approved by the ministerial committee. If the
        specification is now deemed inadequate, the question is why the
        approval process did not catch that, not why the system followed
        its instructions.}
        \item \textit{Agentic systems can be designed with meaningful human
        control built in. It requires investment in review infrastructure
        and tolerance for slower decision cycles. Those are political and
        budgetary choices, not technical constraints.}
    \end{itemize}


\newthought{Key Learnings}
Agency is a spectrum, not a binary: the relevant policy question is not whether AI acts autonomously,
but in which domains, to what degree, and under what oversight. R.U.R.\ and Der Zauberlehrling
together illustrate the alignment problem in miniature: a system optimising for an aggregate goal
may systematically harm those whose interests are underweighted in its objective function; the poem
compresses the arc into fifteen stanzas, the play traces its political consequences across three acts.
Agentic workflows are already deployed in finance, logistics, and digital advertising; the governance
gap they create is a present political problem. Meaningful democratic control over AI agency requires
more than audit trails: legible goals, clear mandates, defined override conditions, and institutions
capable of exercising them. The four themes of this course (Anthropomorphism, Accountability,
Embodiment, Agency) are not independent: an agent that looks human, acts without oversight,
in physical space, with no clear accountability chain, concentrates all four challenges at once.

\medskip
\noindent\textit{Teaser: what governance architecture would you build
to contain such an agent?}

%\input{chapters/05_Maschinenquaelerei.tex}

% Appendix
\appendix

\input{chapters/aa_PreparationGuidance.tex} % → Appendix A
\input{chapters/ac_DebateProtocol.tex}   % → Appendix B


\bibliographystyle{abbrvnat}
\bibliography{philosophyofai_references}



% Hey Mark this is for you, when thinking about where to take this next.
% Think Ethics for Nerds
% Think Book on The Practice of Philosophy a handbook for beginners.
% Think about inviting guest speakers to talk about their work and how it connects to the themes of the course.

\end{document}
